% GENERATED FILE. Edit canonical lessons, then run npm run generate:books.
% canonical-source-hash: fnv1a64:ef8e5db037e8b4af
% canonical-lessons: HI-C41-big, HI-C41-small, HI-C41-good, HI-C41-new, HI-C41-old, HI-C41-adjective-system

\chapter{Describing Things}
\label{ch:hi-describing-things}
\hlchaptermodality{41}

\begin{chapteropening}
\textbf{By the end of this chapter:} \emph{I can point at something in Hindi and ask who or where it is — and I can say what the six words have in common.}

Say the near word, the far word and the question word in a row, and name the part that changed.
\end{chapteropening}

\section[baṛā]{\hi{बड़ा} (baṛā) — big}

\label{lesson:HI-C41-big}



\begin{quote}
You can already name things. This is how you POINT at them.
\end{quote}

\subsection*{You'll want to know: \hi{बड़ा}}
\textbf{\hi{बड़ा}} — \emph{baṛā} — big.

Put it in front of a word you already know and you have said something new. That is all an adjective does in this language.

The last lesson of this chapter gives the one rule that governs every adjective you will ever meet here.

\subsection*{Your turn}
Say these aloud:

\begin{itemize}
\raggedright
  \item \textquotedblleft{}\hi{बड़ा}\textquotedblright{} three times, pointing at something different each time
  \item it once with a word you already know after it
\end{itemize}

\begin{tcolorbox}[breakable,colback=teal!4,colframe=teal!35!black,title={Before you move on}]
What is \emph{baṛā}? (\textbf{big}.) Where does it go — before the word it describes, or after? (\textbf{Before.})
\end{tcolorbox}
\section[choṭā]{\hi{छोटा} (choṭā) — small}

\label{lesson:HI-C41-small}



\begin{quote}
You can already name things. This is how you POINT at them.
\end{quote}

\subsection*{You'll want to know: \hi{छोटा}}
\textbf{\hi{छोटा}} — \emph{choṭā} — small.

Put it in front of a word you already know and you have said something new. That is all an adjective does in this language.

You met \textbf{\hi{बड़ा}} \emph{baṛā} a moment ago. Put them side by side and nothing about either one changes.

The last lesson of this chapter gives the one rule that governs every adjective you will ever meet here.

\subsection*{Your turn}
Say these aloud:

\begin{itemize}
\raggedright
  \item \textquotedblleft{}\hi{छोटा}\textquotedblright{} three times, pointing at something different each time
  \item it once with a word you already know after it
\end{itemize}

\begin{tcolorbox}[breakable,colback=teal!4,colframe=teal!35!black,title={Before you move on}]
What is \emph{choṭā}? (\textbf{small}.) Where does it go — before the word it describes, or after? (\textbf{Before.})
\end{tcolorbox}
\section[acchā]{\hi{अच्छा} (acchā) — good}

\label{lesson:HI-C41-good}



\begin{quote}
You can already name things. This is how you POINT at them.
\end{quote}

\subsection*{You'll want to know: \hi{अच्छा}}
\textbf{\hi{अच्छा}} — \emph{acchā} — good.

Put it in front of a word you already know and you have said something new. That is all an adjective does in this language.

You met \textbf{\hi{छोटा}} \emph{choṭā} a moment ago. Put them side by side and nothing about either one changes.

The last lesson of this chapter gives the one rule that governs every adjective you will ever meet here.

\subsection*{Your turn}
Say these aloud:

\begin{itemize}
\raggedright
  \item \textquotedblleft{}\hi{अच्छा}\textquotedblright{} three times, pointing at something different each time
  \item it once with a word you already know after it
\end{itemize}

\begin{tcolorbox}[breakable,colback=teal!4,colframe=teal!35!black,title={Before you move on}]
What is \emph{acchā}? (\textbf{good}.) Where does it go — before the word it describes, or after? (\textbf{Before.})
\end{tcolorbox}
\section[nayā]{\hi{नया} (nayā) — new}

\label{lesson:HI-C41-new}



\begin{quote}
You can already name things. This is how you POINT at them.
\end{quote}

\subsection*{You'll want to know: \hi{नया}}
\textbf{\hi{नया}} — \emph{nayā} — new.

Put it in front of a word you already know and you have said something new. That is all an adjective does in this language.

You met \textbf{\hi{अच्छा}} \emph{acchā} a moment ago. Put them side by side and nothing about either one changes.

The last lesson of this chapter gives the one rule that governs every adjective you will ever meet here.

\subsection*{Your turn}
Say these aloud:

\begin{itemize}
\raggedright
  \item \textquotedblleft{}\hi{नया}\textquotedblright{} three times, pointing at something different each time
  \item it once with a word you already know after it
\end{itemize}

\begin{tcolorbox}[breakable,colback=teal!4,colframe=teal!35!black,title={Before you move on}]
What is \emph{nayā}? (\textbf{new}.) Where does it go — before the word it describes, or after? (\textbf{Before.})
\end{tcolorbox}
\section[purānā]{\hi{पुराना} (purānā) — old — of a thing, not of a person}

\label{lesson:HI-C41-old}



\begin{quote}
You can already name things. This is how you POINT at them.
\end{quote}

\subsection*{You'll want to know: \hi{पुराना}}
\textbf{\hi{पुराना}} — \emph{purānā} — old — of a thing, not of a person.

Put it in front of a word you already know and you have said something new. That is all an adjective does in this language.

You met \textbf{\hi{नया}} \emph{nayā} a moment ago. Put them side by side and nothing about either one changes.

The last lesson of this chapter gives the one rule that governs every adjective you will ever meet here.

\subsection*{Your turn}
Say these aloud:

\begin{itemize}
\raggedright
  \item \textquotedblleft{}\hi{पुराना}\textquotedblright{} three times, pointing at something different each time
  \item it once with a word you already know after it
\end{itemize}

\begin{tcolorbox}[breakable,colback=teal!4,colframe=teal!35!black,title={Before you move on}]
What is \emph{purānā}? (\textbf{old — of a thing, not of a person}.) Where does it go — before the word it describes, or after? (\textbf{Before.})
\end{tcolorbox}
\section[before the noun, and it AGREES]{before the noun, and it AGREES — five words, one rule}

\label{lesson:HI-C41-adjective-system}



\begin{quote}
Say the five words from this chapter out loud, then say each one in front of a noun you already know.
\end{quote}

\begin{grammarlens}[title={Grammar lens: where an adjective goes}]
You did not learn five separate words. You learned \textbf{where an adjective goes}, and five words that go there.

\noindent
\begin{tabularx}{\linewidth}{@{}>{\raggedright\arraybackslash}X>{\raggedright\arraybackslash}X@{}}
\toprule
\textbf{meaning} & \textbf{word} \\
\midrule
big & \textbf{\hi{बड़ा}} \emph{baṛā} \\
small & \textbf{\hi{छोटा}} \emph{choṭā} \\
good & \textbf{\hi{अच्छा}} \emph{acchā} \\
\bottomrule
\end{tabularx}

Hindi is the one that changes: \hi{बड़ा} for a masculine noun, \hi{बड़ी} for a feminine one.

This is worth more than the five words are. Every noun you already know — your name, the house, the water, the day — can now take any of them in front, and every adjective you meet from here goes in the same place.
\end{grammarlens}

\subsection*{Your turn}
\begin{itemize}
\raggedright
  \item \emph{Say it:} each of the five in front of a noun you already know
  \item \emph{Say it:} two of them with the same noun, one after the other
  \item \emph{Notice:} whether the adjective changed shape when the noun did
\end{itemize}

\begin{tcolorbox}[breakable,colback=teal!4,colframe=teal!35!black,title={Before you move on}]
Where does an adjective go in this language? (\textbf{In front of the noun.}) Does it change shape to match the noun? (and it AGREES.)
\end{tcolorbox}
